\documentclass[aspectratio=43,11pt]{beamer}

%CORE PACKAGES---------------------------------
\usepackage{amsmath, amssymb, graphicx, booktabs, tabularx, setspace, soul}
\usepackage[absolute, overlay]{textpos}
\usepackage{subcaption}
\usepackage{tikz, pgfplots}
\pgfplotsset{compat=1.18}
\usepgfplotslibrary{fillbetween}
\usetikzlibrary{positioning, arrows.meta, arrows.meta, positioning, shapes, calc}
\usepackage[many]{tcolorbox}
\usepackage[ruled,vlined]{algorithm2e}
\usepackage{minted}

%FONTS & STYLING-------------------------------
\usepackage{fontspec, emoji, tikzsymbols}
\usefonttheme{professionalfonts}
\setsansfont[
UprightFont    = *Light,
BoldFont       = *Bold,
ItalicFont     = *Light Italic,
BoldItalicFont = *Bold Italic,
Ligatures      = TeX,
Scale          = MatchLowercase
]{Frutiger LT Pro}
\renewcommand{\familydefault}{\sfdefault}
\linespread{1.}\selectfont
\hyphenpenalty=10000
\exhyphenpenalty=10000
\newcommand{\prob}{\mathsf{P}}
\newcommand{\lik}{\mathcal{L}}
\newcommand{\E}{\mathbb{E}}
\newcommand{\Prb}{\mathsf{P}}
\newcommand{\fixed}[1]{\textcolor{orange}{#1}}
\newcommand{\BetaDist}{\operatorname{Beta}}
\newcommand{\Bin}{\operatorname{Binomial}}
\newcommand{\BB}{\operatorname{BetaBinomial}}
\newcommand{\Bern}{\operatorname{Bernoulli}}
\newcommand{\Var}{\operatorname{Var}}
\newcommand{\given}{\,\mid\,}

%COLORS & THEMES-------------------------------
\definecolor{aggiemaroon}{RGB}{80,0,0}
\definecolor{beamer@blendedblue}{RGB}{80,0,0}
\definecolor{AggieMaroon}{HTML}{500000}

\usecolortheme{default}

\setbeamertemplate{itemize items}{•}
\setbeamertemplate{itemize subitem}[•]
\setbeamertemplate{itemize subsubitem}[•]
\renewcommand{\UrlFont}{\ttfamily\footnotesize}
\newcommand{\vocab}[1]{\textcolor{AggieMaroon}{\textbf{#1}}}
\newcommand{\question}[1]{\begin{block}{Question}#1\end{block}}
\newcommand{\smallnote}[1]{{\footnotesize\textcolor{DarkGray}{#1}}}

%HYPERLINKS------------------------------------
\usepackage{hyperref}
\hypersetup{
colorlinks=true,
linkcolor=aggiemaroon,
urlcolor=aggiemaroon,
citecolor=aggiemaroon
}

%BEAMER TEMPLATES------------------------------
\setbeamertemplate{caption}[numbered]
\setbeamertemplate{navigation symbols}{}
\setbeamertemplate{footline}{
\color{aggiemaroon}{\hfill \insertframenumber/\inserttotalframenumber \hspace{2pt} \vskip2pt}
}

%CUSTOM ENVIRONMENTS---------------------------
\newenvironment{stepenumerate}{
\begin{enumerate}[<+->]
}{\end{enumerate}}

\newenvironment{stepitemize}{
\begin{itemize}[<+->]
}{\end{itemize}}

\newenvironment{stepenumeratewithalert}{
\begin{enumerate}[<+-| alert@+>]
}{\end{enumerate}}

\newenvironment{stepitemizewithalert}{
\begin{itemize}[<+-| alert@+>]
}{\end{itemize}}

%DOCUMENT INFO----------------------------------
\title{Decision and Risk Analysis}
\subtitle{Lecture 3: \textbf{The Beta-Binomial Bayesian Model}}
\author{\emph{}}
\date{\color{aggiemaroon}{
DAEN 427/ISEN 427/ISEN 627\\
\textbf{Alexander Abuabara}\\
Fall 2026
}}
\titlegraphic{\vfill 
\includegraphics[height=1.2cm]{TAMU_logo.png}}

\begin{document}

%-----------------------------------------------
\begin{frame}[plain]
\titlepage
\end{frame}

{
\setbeamercolor{background canvas}{bg=green!10}
%-----------------------------------------------
\begin{frame}{Decision begins before the calculation}
\large
Michelle's historical support is near 45\%.\\
A new poll reports 30 supporters among 50 voters.

\vspace{1em}
\begin{block}{Decision question}
Should the campaign allocate scarce field resources to Michelle's state if winning requires support above 50\%?
\end{block}
\end{frame}
}

%-----------------------------------------------
\begin{frame}{Learning Objectives}
\begin{itemize}
\item Encode prior knowledge about an unknown proportion with a Beta model;
\item Distinguish a Binomial data model from its likelihood function;
\item Update Beta prior parameters after observing Binomial data;
\item Translate a posterior into probabilities relevant to a decision;
\item Choose an action by minimizing posterior expected loss;
\item Test whether the recommendation is robust to prior assumptions.
\end{itemize}
\end{frame}

%-----------------------------------------------
\begin{frame}{Four ingredients organize the analysis}
\begin{center}
\begin{tikzpicture}[node distance=0.75cm, every node/.style={align=center}]
\node[draw,rounded corners,minimum width=2.25cm,minimum height=1.0cm] (p) {Prior\\$\pi\sim\mathrm{Beta}(\alpha,\beta)$};
\node[draw,rounded corners,minimum width=2.25cm,minimum height=1.0cm,right=of p] (d) {Data model\\$Y\mid\pi\sim\mathrm{Bin}(n,\pi)$};
\node[draw,rounded corners,minimum width=2.25cm,minimum height=1.0cm,below=of d] (post) {Posterior\\$\pi\mid y$};
\node[draw,rounded corners,minimum width=2.25cm,minimum height=1.0cm,left=of post] (a) {Action\\minimize expected loss};
\draw[->,thick,black] (p)--(d);
\draw[->,thick,black] (d)--(post);
\draw[->,thick,black] (post)--(a);
\draw[->,thick,black] (p)--(post);
\end{tikzpicture}
\end{center}
\begin{block}{Separation Principle}
The \textbf{probability model} describes uncertainty;\\
the \textbf{loss model} describes preferences and consequences.
\end{block}
\end{frame}

%-----------------------------------------------
\begin{frame}
\begin{block}{Beta-Binomial Conjugacy}
\begin{center}
\begin{tikzpicture}[node distance=0.7cm]
\node[draw, rounded corners, fill=blue!8, minimum width=2.7cm, minimum height=0.7cm] (prior) {Beta prior};
\node[draw, rounded corners, fill=green!10, right=of prior, minimum width=3.0cm, minimum height=0.7cm] (data) {Binomial data};
\node[draw, rounded corners, fill=orange!12, right=of data, minimum width=3.0cm, minimum height=0.7cm] (post) {Beta posterior};
\draw[-{Latex[length=2mm]}, thick] (prior) -- (data);
\draw[-{Latex[length=2mm]}, thick] (data) -- (post);
\end{tikzpicture}
\end{center}
\end{block}
\end{frame}

%-----------------------------------------------
\begin{frame}{Conjugate}
\begin{itemize}
\item \textbf{General definition:} joined together, coupled, or connected (especially in pairs).
\item \textbf{Grammar and linguistics:} changing a verb's form to show tense, person, number, or mood; for example, \emph{run}, \emph{runs}, \emph{ran}, and \emph{running}.
\item \textbf{Mathematics:} paired expressions with a related form.
\begin{itemize}
\item The conjugate of $x+3$ is $x-3$.
\item The conjugate of $a+bi$ is $a-bi$.
\end{itemize}
\item \textbf{Chemistry:} linked substances, including acid--base pairs related by proton gain or loss.
\item \textbf{Biology:} temporary pairing or fusion of single-celled organisms to exchange genetic material.
\end{itemize}
\end{frame}

%-----------------------------------------------
\begin{frame}{}
In Bayesian statistics, conjugate refers to a special mathematical relationship between a \ul{prior probability distribution} and a \ul{likelihood function}.
\begin{itemize}
\item When a prior distribution is combined with a likelihood to calculate the posterior distribution, it is called a \textbf{conjugate prior} \ul{if the resulting posterior belongs to the exact same probability distribution family as the prior}.
\end{itemize}
\end{frame}

%-----------------------------------------------
\begin{frame}{}
\begin{itemize}
\item \textbf{Simplifies Math}\\Bayesian updates usually require complex calculus (integration) to solve.
With a conjugate prior, the math simplifies into basic algebra. You can calculate the posterior distribution by simply updating the parameters of the prior distribution. {\it (You will see!)}
\item \textbf{Saves Computing Power}\\Because bypasses intense calculations, allows for instant updating without needing heavy simulation algorithms (like Markov Chain Monte Carlo, MCMC).
\end{itemize}
\end{frame}

%-----------------------------------------------
\begin{frame}{}
In practice, certain likelihoods pair perfectly with specific priors:
\begin{itemize}
\item \textbf{Beta-Binomial Conjugacy:} Used for binary (yes/no) outcomes like coin flips or click-through rates. If your prior is a Beta distribution and your data is Binomial, your posterior will also be a Beta distribution.
\item \textbf{Normal-Normal Conjugacy:} Used for continuous data (like heights or temperatures). If your prior is a Normal (Gaussian) distribution and your data is normally distributed, your posterior will also be a Normal distribution.
\item \textbf{Dirichlet-Multinomial Conjugacy:} Used for outcomes with multiple categories (like rolling a 6-sided die or text categorization). It is the multi-category extension of the Beta-Binomial pair.
\end{itemize}
\end{frame}

%-----------------------------------------------
\begin{frame}{}
To understand conjugate priors mathematically, we use Bayes rules.
It states that the posterior probability is proportional to the prior multiplied by the likelihood:
\[
\text{Posterior} \propto \text{Likelihood}\times\text{Prior}
\]
\[
P(\theta\mid x)\propto P(x\mid\theta)P(\theta)
\]

\pause
\vspace{.5cm}
This models binary outcomes, for instance: predicting a coin's true probability of landing heads $(e)$ after seeing a series of flips $(x)$.
\end{frame}

{
\setbeamercolor{background canvas}{bg=gray!10}
%-----------------------------------------------
\begin{frame}{The Beta distribution ...}
\footnotesize
\begin{description}
\item[1730s] \textbf{Leonhard Euler} investigated the integral behind
the Beta function while generalizing factorials.

\item[1763] \textbf{Thomas Bayes} studied \emph{inverse probability}:
inferring an unknown Bernoulli success probability from observed trials.
His posterior had the structure now associated with a Beta distribution.

\item[1839] \textbf{Jacques Binet} formalized the integral and used the
notation $B$:
\[
B(\alpha,\beta)
=\int_0^1 x^{\alpha-1}(1-x)^{\beta-1}\,dx.
\]

\item[1895] \textbf{Karl Pearson} included the curve in his system as the
\emph{Pearson~Type~I~distribution}, useful for bounded, skewed data.
\end{description}

\begin{alertblock}{}
By the mid-20th century, \emph{Beta distribution} became the standard name.
\end{alertblock}
\end{frame}

%-----------------------------------------------
\begin{frame}{The Beta distribution ... intuition}
\footnotesize
The Beta numerator is very similar to the Binomial distribution pmf:
\[
f(x)
=
\frac{x^{\alpha-1}(1-x)^{\beta-1}}{B(\alpha,\beta)}
\qquad \longrightarrow \qquad
\text{Binomial pmf}
=
\binom{n}{x}p^x(1-p)^{n-x}.
\]

The Binomial and Beta distributions are very similar apart from \ul{one key thing}:

\begin{itemize}
\item For the {\bf Beta distribution}, \hl{the probability is a random variable} that we are trying to estimate.

\item For the {\bf Binomial distribution}, \hl{the probability is a fixed parameter} that we use to infer the probability of \(n\) successes.
\end{itemize}

Therefore, we can use the Beta distribution to estimate the probability of an event if we know the number of {\it successes} ($\alpha + 1$) and the number of {\it failures} ($\beta - 1$).
\end{frame}

%-----------------------------------------------
\begin{frame}{The Beta distribution ... flexible and easy to update}
\footnotesize
\begin{columns}[T,onlytextwidth]
\begin{column}{0.46\textwidth}
\vspace{1cm}
The Beta distribution is a family of continuous probability distributions defined on $[0,1]$:
\[
f(x\mid\alpha,\beta)
=\frac{x^{\alpha-1}(1-x)^{\beta-1}}
 {B(\alpha,\beta)}
\]
The parameters ($\alpha,\beta>0$) produce symmetric, skewed, U-shaped, or nearly uniform densities on $[0,1]$.
This makes it well suited for percentages, proportions, and probabilities.
\end{column}

\begin{column}{0.48\textwidth}
\begin{block}{Applications}
\begin{itemize}
\item \textbf{A/B Testing:} conversion / click-through rates.
\item \textbf{Project Management (PERT Analysis):} uncertain task durations from optimistic, likely, and pessimistic estimates.
\item \textbf{Glicko Rating System:} uncertainty about an unknown probability, used in competitive chess and video games to model a player's underlying skill level along with the system's current uncertainty about that skill.
\item \textbf{Reliability:} failure or success rates updated as evidence arrives.
\end{itemize}
\end{block}
\end{column}
\end{columns}
\end{frame}
}

%-----------------------------------------------
\begin{frame}{Step 1: Define the Prior $P(\theta)$}
We model our initial belief about $\theta$ using a Beta distribution.

It is defined by two shape parameters, $\alpha$ and $\beta$, which act like "pseudo-counts" for successes and failures:
\[
\theta\sim\operatorname{Beta}(\alpha,\beta),\qquad
p(\theta)=\frac{1}{B(\alpha,\beta)}\theta^{\alpha-1}(1-\theta)^{\beta-1}
\]

\pause
\vspace{.75cm}
Since the denominator $B(\alpha,\beta)$ is just a normalizing constant (a fixed number to make the total probability equal 1), we can ignore it during the calculation and focus on the core variables:
\[
p(\theta)\propto\theta^{\alpha-1}(1-\theta)^{\beta-1}
\]
\end{frame}

%-----------------------------------------------
\begin{frame}{Step 2: Define the Likelihood $P(x\mid\theta)$}
Suppose a coin is flipped $n$ times and $k$ heads are observed.
The probability of getting this exact result follows a Binomial distribution:
\[
P(x\mid\theta)=\binom{n}{k}\theta^k(1-\theta)^{n-k}
\]

\pause
\vspace{.75cm}
Again, the combination formula $\binom{n}{k}$ is a constant relative to $\theta$, so we drop it for the proportional calculation:
\[
P(x\mid\theta)\propto\theta^k(1-\theta)^{n-k}
\]
\end{frame}

%-----------------------------------------------
\begin{frame}{Steps 3: Multiply Likelihood by Prior}
Multiply the core terms of the likelihood and the prior to find the shape of the posterior:
\begin{align*}
P(\theta\mid x)
&\propto P(x\mid\theta)P(\theta)\\
&\propto [\theta^k(1-\theta)^{n-k}]
[\theta^{\alpha-1}(1-\theta)^{\beta-1}]
\end{align*}
\end{frame}

%-----------------------------------------------
\begin{frame}{Steps 4: Combine Exponents}
Multiply likelihood and prior:
\begin{align*}
P(\theta\mid x)
&\propto P(x\mid\theta)P(\theta)\\
&\propto [\theta^k(1-\theta)^{n-k}]
[\theta^{\alpha-1}(1-\theta)^{\beta-1}]\\
&\propto \theta^{k+\alpha-1}(1-\theta)^{n-k+\beta-1}\\
&=\theta^{(\alpha+k)-1}(1-\theta)^{[\beta+(n-k)]-1}
\end{align*}
\tiny Ps. Remember the algebraic rule $x^a x^b=x^{a+b}$
\end{frame}

%-----------------------------------------------
\begin{frame}{Step 5: Identify the Posterior}
The result has the same form as the original Beta distribution, but with updated parameters:
\[
\theta\mid x\sim\operatorname{Beta}\bigl(\alpha+k,\;\beta+n-k\bigr).
\]
\begin{center}
\begin{tabular}{>{\bfseries}lll}
\toprule
Parameter &  & Posterior update\\
\midrule
$\alpha$ & {\footnotesize \it Successes (heads)} & $\alpha_{\text{post}}=\alpha_{\text{prior}}+k$\\
$\beta$  & {\footnotesize \it Failures (tails)}  & $\beta_{\text{post}}=\beta_{\text{prior}}+(n-k)$\\
\bottomrule
\end{tabular}
\end{center}
\end{frame}

%-----------------------------------------------
\begin{frame}{Example: Testing a Coin}
\small
Let's say you want to estimate the true probability $\theta$ that a coin lands on heads.
\begin{enumerate}
\item \textbf{Prior:} choose $\operatorname{Beta}(2,2)$, representing a slightly uncertain belief in a fair coin.
This acts like you've already seen 2 heads ($\alpha$) and 2 tails ($\beta$).
\item \textbf{Data:} flip the coin $n=10$ times and observe $k=7$ heads / $3$ tails.
\item \textbf{Posterior update:}
\[
\alpha_{\text{post}}=2+7=9,
\qquad
\beta_{\text{post}}=2+3=5.
\]
\end{enumerate}
Your updated belief about the coin is now captured by a $\operatorname{Beta}(9,5)$ distribution!
\end{frame}

%-----------------------------------------------
\begin{frame}[fragile]{}
\begin{minted}[fontsize=\scriptsize, linenos]{r}
# 1. Define prior parameters: Beta(2, 2)
alpha_prior <- 2
beta_prior  <- 2

# 2. Enter observed data: 7 heads in 10 flips
n_flips <- 10
heads <- 7
tails <- n_flips - heads

# 3. Compute posterior parameters: Beta(9, 5)
alpha_posterior <- alpha_prior + heads
beta_posterior  <- beta_prior + tails

# 4. Create theta values
x <- seq(0, 1, length.out = 500)
\end{minted}
\end{frame}

%-----------------------------------------------
\begin{frame}[fragile]{}
\begin{minted}[fontsize=\scriptsize, linenos]{r}
# Prior density
prior_curve <- dbeta(x,
                     shape1 = alpha_prior,
                     shape2 = beta_prior)

# Binomial likelihood as a function of theta
likelihood_curve <- dbinom(heads,
                           size = n_flips,
                           prob = x)

# Normalize likelihood for comparison with the densities
dx <- x[2] - x[1]
likelihood_curve <- likelihood_curve / sum(likelihood_curve * dx)

# Posterior density
posterior_curve <- dbeta(x,
                         shape1 = alpha_posterior,
                         shape2 = beta_posterior)
\end{minted}
\end{frame}

%-----------------------------------------------
\begin{frame}[fragile]{}
\begin{minted}[fontsize=\scriptsize, linenos]{r}
# 5. Plot all three curves
y_max <- max(prior_curve, likelihood_curve, posterior_curve)

plot(x, prior_curve,
     type = "l", col = "red",
     lwd = 3, lty = 2,
     ylim = c(0, y_max),
     main = "Bayesian Update",
     xlab = expression("Probability of Heads (" * theta * ")"),
     ylab = "Normalized Density")

lines(x, likelihood_curve, col = "darkgreen", lwd = 3, lty = 3)

lines(x, posterior_curve, col = "blue", lwd = 3, lty = 1)

legend("topleft",
  legend = c(
    "Prior: Beta(2, 2)",
    "Likelihood: 7 heads in 10 flips",
    "Posterior: Beta(9, 5)"),
  col = c("red", "darkgreen", "blue"), lty = c(2, 3, 1), lwd = 3, bty = "n")
\end{minted}
\end{frame}

%-----------------------------------------------
\begin{frame}{}
\begin{center}
\includegraphics[height=.9\textheight]{plot_1_beta_binomial.png}
\end{center}
\end{frame}

%-----------------------------------------------
\begin{frame}{}
\begin{itemize}
\item The red prior is low and spread out, representing initial uncertainty.
\item The blue posterior shifts right and becomes taller and narrower.
\item Its mean is $9/(9+5)\approx0.64$, indicating support for a heads-biased coin.
\end{itemize}
\end{frame}

%-----------------------------------------------
\begin{frame}{What If Prior Knowledge Is a Proportion?}
\small
Suppose prior knowledge says the success rate is about $60\%$. A single value such as $0.60$ is not a distribution; by itself it implies absolute certainty.

Represent both the center and the strength of the belief with a Beta distribution. Its expected value is
\[
\mu=E(\theta)=\frac{\alpha}{\alpha+\beta}.
\]

Choose a pseudo-sample size $N$ to express certainty:
\begin{itemize}
\item A small $N$ means the rate is believed to be near $60\%$, but with high uncertainty.
\item A large $N$ expresses strong confidence, as if many prior trials had already been observed.
\end{itemize}
\end{frame}

%-----------------------------------------------
\begin{frame}{Convert a Proportion to Beta Parameters}
For expected proportion $\mu$ and prior certainty $N$,
\[
\alpha_{\text{prior}}=\mu N,
\qquad
\beta_{\text{prior}}=N-\alpha_{\text{prior}}=(1-\mu)N.
\]

\textbf{For $\mu=0.60$:}
\begin{itemize}
\item \textbf{Weak prior, $N=5$}\\$\alpha=0.60(5)=3$, $\beta=5-3=2$; hence $\operatorname{Beta}(3,2)$.
\item \textbf{Strong prior, $N=50$}\\$\alpha=0.60(50)=30$, $\beta=50-30=20$; $\operatorname{Beta}(30,20)$.
\end{itemize}
\end{frame}

%-----------------------------------------------
\begin{frame}{}
\begin{center}
\includegraphics[height=.9\textheight]{plot_2_beta_beta.png}
\end{center}
\end{frame}

%-----------------------------------------------
\begin{frame}{}
\footnotesize
The plot compares the probability density functions of two Beta distributions for the parameter $\theta$:
$\operatorname{Beta}(3,2)$ and $\operatorname{Beta}(30,20)$.

For a Beta distribution with parameters $\alpha$ and $\beta$, the mean is

\[
\mathbb{E}[\theta]
=
\frac{\alpha}{\alpha+\beta}.
\]

Therefore,

\[
\mathbb{E}[\theta \mid \theta \sim \operatorname{Beta}(3,2)]
=
\frac{3}{3+2}
=
0.60,
\]

and

\[
\mathbb{E}[\theta \mid \theta \sim \operatorname{Beta}(30,20)]
=
\frac{30}{30+20}
=
0.60.
\]

Thus, both distributions have the same mean. This explains why the two
vertical mean lines coincide at $\theta=0.60$.
\end{frame}

%-----------------------------------------------
\begin{frame}{}
\footnotesize
The principal difference is their degree of concentration.

\begin{itemize}
\item The $\operatorname{Beta}(3,2)$ distribution is broad and
assigns substantial probability density to a relatively wide range of
values of $\theta$. Represents greater uncertainty about $\theta$.
\item The $\operatorname{Beta}(30,20)$ distribution is narrow and sharply concentrated near $0.60$, representing substantially greater certainty.\\
\end{itemize}

The variance of a Beta distribution is

\[
\operatorname{Var}(\theta)
=
\frac{\alpha\beta}
     {(\alpha+\beta)^2(\alpha+\beta+1)}
\]

\[
\operatorname{Var}(\theta \mid \theta \sim \operatorname{Beta}(3,2))
=
\frac{3 \cdot 2}{(3+2)^2(3+2+1)}
=
0.04
\]

\[
\operatorname{Var}(\theta \mid
\theta \sim \operatorname{Beta}(30,20))
=
\frac{30 \cdot 20}{(30+20)^2(30+20+1)}
\approx
0.0047
\]

The variance of $\operatorname{Beta}(30,20)$ is approximately one-ninth that of 
$\operatorname{Beta}(3,2)$!
\end{frame}

%-----------------------------------------------
\begin{frame}{}
\footnotesize
The highest point of a distribution is its mode, which does not
necessarily coincide with its mean. When $\alpha>1$ and $\beta>1$, the
mode of a Beta distribution is

\[
\operatorname{Mode}(\theta)
=
\frac{\alpha-1}{\alpha+\beta-2}
\]

\[
\operatorname{Mode}(\theta \mid
\theta \sim \operatorname{Beta}(3,2))
=
\frac{3-1}{3+2-2}
=
\frac{2}{3}
\approx
0.667
\]

\[
\operatorname{Mode}(\theta \mid
\theta \sim \operatorname{Beta}(30,20))
=
\frac{30-1}{30+20-2}
=
\frac{29}{48}
\approx
0.604
\]

The blue curve reaches its maximum to the right of its mean,
the peak of the red curve is nearly aligned with its mean.

From a Bayesian perspective, the two distributions express the same
central prior belief, namely that $\theta$ is approximately $0.60$.
However, $\operatorname{Beta}(30,20)$ represents much stronger prior
information (or greater confidence in that belief) than
$\operatorname{Beta}(3,2)$.
\end{frame}

%-----------------------------------------------
\begin{frame}{Problem}
\begin{textblock*}{2cm}(11cm,0.3cm)

\includegraphics[width=1.5cm]{bookmark.png}
\end{textblock*}

\begin{block}{Unknown Success Probability}
We observe binary outcomes, such as success/failure, yes/no, heads/tails.
\end{block}

\vspace{0.2cm}
Let $\theta$ be the unknown probability of success.

\begin{align*}
Y_i \given \theta &\sim \Bern(\theta), \qquad i=1,\dots,n, \\
X = \sum_{i=1}^n Y_i \given \theta &\sim \Bin(n,\theta).
\end{align*}

\begin{block}{}
Use observed successes $X=x$ out of $n$ trials to update beliefs about $\theta$.
\end{block}
\end{frame}

%-----------------------------------------------
\begin{frame}{Prior: the Beta distribution}
A natural prior for a probability $\theta \in (0,1)$ is
\[
\theta \sim \BetaDist(\alpha,\beta), \qquad \alpha>0,\; \beta>0.
\]

Its density is
\[
p(\theta)=\frac{\Gamma(\alpha+\beta)}{\Gamma(\alpha)\Gamma(\beta)}\,\theta^{\alpha-1}(1-\theta)^{\beta-1}.
\]

For conjugacy, we focus on the kernel:
\[
p(\theta) \propto \theta^{\alpha-1}(1-\theta)^{\beta-1}.
\]

\begin{block}{Interpretation}
$\alpha-1$ acts like prior successes; $\beta-1$ acts like prior failures.
\end{block}
\end{frame}

%-----------------------------------------------
\begin{frame}{Likelihood: the Binomial model}
If $X=x$ successes are observed in $n$ independent trials, then
\[
X \given \theta \sim \Bin(n,\theta).
\]

The likelihood is
\[
p(x \given \theta)=\binom{n}{x}\theta^x(1-\theta)^{n-x}.
\]

As a function of $\theta$, constants not involving $\theta$ can be dropped:
\[
L(\theta \given x) \propto \theta^x(1-\theta)^{n-x}.
\]

\begin{block}{Data contribution}
The data contribute $x$ successes and $n-x$ failures.
\end{block}
\end{frame}

%-----------------------------------------------
\begin{frame}{Bayes' rule: multiply prior and likelihood}
Bayes' rule says
\[
p(\theta \given x)=\frac{p(x \given \theta)p(\theta)}{p(x)}
\quad \Longrightarrow \quad
p(\theta \given x) \propto p(x \given \theta)p(\theta).
\]

Substitute the likelihood and prior kernels:
\begin{align*}
p(\theta \given x)
&\propto \Big[\theta^x(1-\theta)^{n-x}\Big]
          \Big[\theta^{\alpha-1}(1-\theta)^{\beta-1}\Big] \\
&= \theta^{x+\alpha-1}(1-\theta)^{n-x+\beta-1}.
\end{align*}

\begin{block}{}
The posterior has the same functional form as the prior.
\end{block}
\end{frame}

%-----------------------------------------------
\begin{frame}{Recognize the posterior distribution}
Compare the posterior kernel
\[
p(\theta \given x) \propto
\theta^{x+\alpha-1}(1-\theta)^{n-x+\beta-1}
\]
with the Beta kernel
\[
\BetaDist(a,b): \qquad \theta^{a-1}(1-\theta)^{b-1}.
\]

Therefore,
\[
\boxed{\theta \given X=x \sim \BetaDist(\alpha+x,\; \beta+n-x)}.
\]

\begin{center}
\begin{tabular}{lll}
\toprule
Quantity & Prior & Posterior \\
\midrule
Success parameter & $\alpha$ & $\alpha+x$ \\
Failure parameter & $\beta$ & $\beta+n-x$ \\
\bottomrule
\end{tabular}
\end{center}
\end{frame}

%-----------------------------------------------
\begin{frame}{Why is called conjugacy?}
\begin{block}{Definition}
A prior is \alert{conjugate} to a likelihood when the posterior belongs to the same family as the prior.
\end{block}

For the beta-binomial case:
\[
\underbrace{\BetaDist(\alpha,\beta)}_{\text{prior for }\theta}
\quad + \quad
\underbrace{\Bin(n,\theta)}_{\text{likelihood for }x}
\quad \Longrightarrow \quad
\underbrace{\BetaDist(\alpha+x,\beta+n-x)}_{\text{posterior for }\theta}.
\]

\vspace{0.25cm}
This works because both prior and likelihood are powers of
\[
\theta \quad \text{and} \quad 1-\theta.
\]
\end{frame}

%-----------------------------------------------
\begin{frame}{Step-by-step numeric example}
\small
Suppose a prior belief is
\[
\theta \sim \BetaDist(2,3),
\]
and we observe $x=7$ successes in $n=10$ trials.

\begin{align*}
\alpha' &= \alpha+x = 2+7 = 9\\
\beta'  &= \beta+n-x = 3+10-7 = 6
\end{align*}

So the posterior is
\[
\boxed{\theta \given X=7 \sim \BetaDist(9,6)}
\]

Posterior mean:
\[
\E[\theta \given x] = \frac{\alpha+x}{\alpha+\beta+n}
= \frac{9}{15}=0.60
\]
\end{frame}

%-----------------------------------------------
\begin{frame}{Posterior predictive: the beta-binomial distribution}
\small
For future data $\widetilde{X}$ from $m$ new trials, integrate over uncertainty in $\theta$:
\[
p(\widetilde{x} \given x)
= \int_0^1 p(\widetilde{x} \given \theta) p(\theta \given x)\,d\theta
\]

Using $\theta \given x \sim \BetaDist(\alpha+x,\beta+n-x)$:
\[
\boxed{
\widetilde{X}\given x \sim \BB\big(m,\alpha+x,\beta+n-x\big)
}
\]

with probability mass function
\[
p(\widetilde{x}\given x)
= \binom{m}{\widetilde{x}}
\frac{B(\widetilde{x}+\alpha+x,\; m-\widetilde{x}+\beta+n-x)}
     {B(\alpha+x,\; \beta+n-x)}
\]
\end{frame}

%-----------------------------------------------
\begin{frame}{Complete derivation in one line}
\begin{align*}
p(\theta \given x)
&\propto p(x\given\theta)p(\theta) \\
&\propto \theta^x(1-\theta)^{n-x}\,\theta^{\alpha-1}(1-\theta)^{\beta-1} \\
&= \theta^{\alpha+x-1}(1-\theta)^{\beta+n-x-1} \\
&\Rightarrow \theta \given x \sim \BetaDist(\alpha+x,\beta+n-x).
\end{align*}

\begin{block}{}
Posterior parameters = prior counts + observed counts.
\[
(\alpha,\beta) \mapsto (\alpha + \text{successes},\; \beta + \text{failures}).
\]
\end{block}
\end{frame}

%-----------------------------------------------
%-----------------------------------------------
%-----------------------------------------------
\begin{frame}{A Beta prior separates location from strength}
For $0<\pi<1$,
\[
f(\pi)=\frac{\Gamma(\alpha+\beta)}{\Gamma(\alpha)\Gamma(\beta)}
\pi^{\alpha-1}(1-\pi)^{\beta-1}
\]
\begin{columns}[T]
\column{0.48\textwidth}
\textbf{Prior center}
\[
\E(\pi)=\frac{\alpha}{\alpha+\beta}
\]
\column{0.48\textwidth}
\textbf{Prior concentration}
\[
\kappa=\alpha+\beta
\]
Larger $\kappa$ means less prior dispersion at a fixed mean.
\end{columns}
\end{frame}

%-----------------------------------------------
\begin{frame}{Equal means can hide very different prior commitments}
\begin{center}
\begin{tikzpicture}
\begin{axis}[width=.88\textwidth,height=.55\textheight,xmin=0,xmax=1,ymin=0,
axis lines=left,xlabel={$\pi$},ylabel={density},ytick=\empty,
legend style={draw=none,at={(.5,-.27)},anchor=north,legend columns=3}]
\addplot[red,thick,domain=.001:.999,samples=150] {6*x*(1-x)};
\addlegendentry{Beta(2,2)}
\addplot[black,thick,dashed,domain=.001:.999,samples=150] {630*x^4*(1-x)^4};
\addlegendentry{Beta(5,5)}
\addplot[blue,thick,dotted,domain=.001:.999,samples=150] {7.976*x^19*(1-x)^19/((0.5)^38)};
\addlegendentry{Beta(20,20)} % , schematic

\end{axis}
\end{tikzpicture}
\end{center}
All three center at 0.5; their effective prior information differs sharply.
\end{frame}

{
\setbeamercolor{background canvas}{bg=green!10}
%-----------------------------------------------
\begin{frame}{The historical polls suggest Beta(45,55)}
\begin{textblock*}{2cm}(11cm,0.3cm)

\includegraphics[width=1.5cm]{bookmark.png}
\end{textblock*}
\begin{columns}[T]
\column{0.52\textwidth}
\[
\pi\sim\mathrm{Beta}(45,55)
\]
\begin{align*}
\E(\pi)&=0.45,\\
\operatorname{Mode}(\pi)&\approx0.449,\\
\operatorname{SD}(\pi)&\approx0.0495.
\end{align*}
\column{0.44\textwidth}
\begin{block}{Interpretation}
The prior is centered below the winning threshold and behaves like roughly 100 units of prior concentration.
\end{block}
\end{columns}
\vspace{0.5em}
The parameters are assumptions to defend, not counts that were literally observed.
\end{frame}

%-----------------------------------------------
\begin{frame}{Check the Binomial assumptions before updating}
The data model is
\[
Y\mid\pi\sim\mathrm{Bin}(n,\pi).
\]
It requires:
\begin{itemize}
\item a fixed number $n$ of trials;
\item a binary outcome under a clear definition of ``success'';
\item a common success probability $\pi$;
\item conditional independence across trials.
\end{itemize}
\begin{block}{Risk}
Sampling bias, nonresponse, clustering, and changing support are not repaired by conjugate algebra.
\end{block}
\end{frame}

%-----------------------------------------------
\begin{frame}{A likelihood ranks parameter values;\\
it is not a probability distribution over them}
After observing $y$ successes in $n$ trials,
\[
L(\pi\mid y)=\binom{n}{y}\pi^y(1-\pi)^{n-y}.
\]
For $y=30$, $n=50$:
\begin{itemize}
\item the likelihood is maximized at $\hat\pi=y/n=0.60$;
\item $L(0.60\mid30)\approx0.115$;
\item $L(0.50\mid30)\approx0.042$.
\end{itemize}
\begin{alertblock}{Common error!!}
$L(0.60\mid30)$ is the probability of the observed count if $\pi=0.60$; it is not $\Prb(\pi=0.60\mid y)$.
\end{alertblock}
\end{frame}
}

%-----------------------------------------------
\begin{frame}{Conjugacy turns updating into addition}
\[
\underbrace{\pi\sim\mathrm{Beta}(\alpha,\beta)}_{\text{prior}},\qquad
\underbrace{Y\mid\pi\sim\mathrm{Bin}(n,\pi)}_{\text{data model}}
\]
\[
\boxed{\pi\mid(Y=y)\sim\mathrm{Beta}(\alpha+y,\;\beta+n-y)}
\]
\begin{center}
\begin{tabular}{lcc}
\toprule
 & Success parameter & Failure parameter\\
\midrule
Prior & $\alpha$ & $\beta$\\
Data & $y$ & $n-y$\\
Posterior & $\alpha+y$ & $\beta+n-y$\\
\bottomrule
\end{tabular}
\end{center}
\end{frame}

%-----------------------------------------------
\begin{frame}{}
\begin{align*}
f(\pi\mid y)
&\propto f(\pi)L(\pi\mid y)\\
&\propto \pi^{\alpha-1}(1-\pi)^{\beta-1}
          \pi^{y}(1-\pi)^{n-y}\\
&\propto \pi^{(\alpha+y)-1}(1-\pi)^{(\beta+n-y)-1}.
\end{align*}
\begin{block}{}
Constants not involving $\pi$ determine normalization,\\
but not the posterior family or its shape.
\end{block}
\end{frame}

{
\setbeamercolor{background canvas}{bg=green!10}
%-----------------------------------------------
\begin{frame}{The poll moves the mean up and uncertainty down}
\[
\mathrm{Beta}(45,55)+\{30\text{ successes},20\text{ failures}\}
\Rightarrow\mathrm{Beta}(75,75).
\]
\begin{center}
\begin{tabular}{lcc}
\toprule
 & Prior & Posterior\\
\midrule
Mean & 0.450 & 0.500\\
Mode & 0.449 & 0.500\\
Standard deviation & 0.0495 & 0.0407\\
\bottomrule
\end{tabular}
\end{center}
\vspace{0.7em}
The estimate shifts toward the new sample proportion, while combined information narrows the distribution.
\end{frame}

%-----------------------------------------------
\begin{frame}{The posterior mean is a credibility-weighted average}
\[
\E(\pi\mid y)
=\frac{\alpha+\beta}{\alpha+\beta+n}\frac{\alpha}{\alpha+\beta}
+\frac{n}{\alpha+\beta+n}\frac{y}{n}.
\]
For the campaign:
\[
0.50=\underbrace{\frac{100}{150}}_{2/3}(0.45)
+\underbrace{\frac{50}{150}}_{1/3}(0.60).
\]
\begin{block}{}
The prior dominates the average because its concentration is twice the new sample size.
\end{block}
\end{frame}

%-----------------------------------------------
\begin{frame}{Posterior probability answers the operational question}
Suppose ``viable'' means $\pi>0.50$.

Then
\[
q=\Prb(\pi>0.50\mid y)
=1-F_{\mathrm{Beta}(75,75)}(0.50)=0.50.
\]
The posterior is symmetric, so the evidence leaves equal probability above and below the threshold.
\begin{block}{But this is not yet a decision}
A 50\% chance of winning can justify investing or withdrawing, depending on the losses attached to each error.
\end{block}
\end{frame}

%-----------------------------------------------
\begin{frame}{A loss table converts belief into action}
Let $a_I$ denote ``invest'' and $a_W$ denote ``withdraw.''
\begin{center}
\begin{tabular}{lcc}
\toprule
 & $\pi>0.50$ & $\pi\le0.50$\\
\midrule
Invest $a_I$ & $0$ & $C_{FP}$\\
Withdraw $a_W$ & $C_{FN}$ & $0$\\
\bottomrule
\end{tabular}
\end{center}
With $q=\Prb(\pi>0.50\mid y)$,
\[
\rho(a_I)=C_{FP}(1-q),\qquad \rho(a_W)=C_{FN}q.
\]
Choose invest when
\[
\boxed{q>\frac{C_{FP}}{C_{FP}+C_{FN}}}.
\]

\tiny
\begin{itemize}
\item $C_{FP}$ (false-positive cost): The cost of investing when Michelle’s true support is not above 50\%.
\item $C_{FN}$ (false-negative cost): The cost of withdrawing when her true support actually is above 50\%.
\end{itemize}
\end{frame}

%-----------------------------------------------
\begin{frame}{The optimal action depends on the cost ratio}
Here $q=0.50$.
\begin{center}
\begin{tabular}{lccc}
\toprule
$C_{FP}$ & $C_{FN}$ & Invest threshold & Recommendation\\
\midrule
1 & 1 & 0.50 & indifferent\\
1 & 3 & 0.25 & invest\\
3 & 1 & 0.75 & withdraw\\
\bottomrule
\end{tabular}
\end{center}
\begin{block}{}
If missing a winnable state is three times as costly as investing in a losing state, the same posterior supports investment.
\end{block}
\end{frame}

%-----------------------------------------------
\begin{frame}{Sensitivity analysis asks whether the decision is fragile}
Keep the prior mean at 0.45 but vary its concentration $\kappa$:
\[
\alpha=0.45\kappa,\qquad\beta=0.55\kappa.
\]
\begin{center}
\begin{tabular}{rccc}
\toprule
$\kappa$ & Prior & Posterior mean & Approx. $\Prb(\pi>0.5\mid y)$\\
\midrule
20 & Beta(9,11) & 0.557 & 0.83\\
60 & Beta(27,33) & 0.518 & 0.65\\
100 & Beta(45,55) & 0.500 & 0.50\\
\bottomrule
\end{tabular}
\end{center}
\begin{block}{Implication}
A recommendation near its loss-based threshold should trigger prior and model sensitivity checks.
\end{block}
\end{frame}
}

%-----------------------------------------------
\begin{frame}{Posterior predictive risk concerns what happens next}
For a future sample of size $m$,
\[
\tilde Y\mid y\sim\mathrm{BetaBinomial}
\bigl(m,\alpha+y,\beta+n-y\bigr).
\]
\[
\Prb(\tilde Y=k\mid y)=\binom{m}{k}
\frac{B(\alpha+y+k,\,\beta+n-y+m-k)}
{B(\alpha+y,\,\beta+n-y)}.
\]
\begin{block}{Use the right uncertainty}
$\pi\mid y$ describes the latent support rate; $\tilde Y\mid y$ describes a future observed count and includes sampling variability.
\end{block}
\end{frame}

%-----------------------------------------------
\begin{frame}[fragile]{A compact R workflow}
\begin{minted}[fontsize=\scriptsize, linenos]{r}
alpha <- 45; beta <- 55
y <- 30; n <- 50

a_post <- alpha + y
b_post <- beta + n - y

post_mean <- a_post / (a_post + b_post)
p_viable  <- 1 - pbeta(.50, a_post, b_post)
ci_95     <- qbeta(c(.025, .975), a_post, b_post)

# posterior predictive simulation for m = 100
pi_draw <- rbeta(100000, a_post, b_post)
y_new   <- rbinom(100000, 100, pi_draw)
\end{minted}
\end{frame}

%-----------------------------------------------
\begin{frame}{Simulation gives a second route to posterior intuition}
\begin{enumerate}
\item Draw $\pi^{(s)}\sim\mathrm{Beta}(45,55)$.
\item Draw $Y^{(s)}\sim\mathrm{Bin}(50,\pi^{(s)})$.
\item Retain draws for which $Y^{(s)}=30$.
\item The retained $\pi^{(s)}$ values approximate $\mathrm{Beta}(75,75)$.
\end{enumerate}
\begin{block}{}
The algebra is exact and efficient; simulation makes conditioning concrete and extends to models without closed-form updates.
\end{block}
\end{frame}

%-----------------------------------------------
\begin{frame}{Model risk can dominate parameter uncertainty}
\footnotesize
\begin{center}
\begin{tabular}{
>{\raggedright\arraybackslash}p{0.27\textwidth}
>{\raggedright\arraybackslash}p{0.28\textwidth}
>{\raggedright\arraybackslash}p{0.34\textwidth}
}
\toprule
Assumption & Failure mode & Response \\
\midrule
Representative sample & Selection/nonresponse bias & Weighting, redesign, bias analysis \\ \hline
Conditional independence & Clustered responses & Hierarchical or overdispersed model \\ \hline
Constant $\pi$ & Time-varying support & Dynamic model; recency weighting \\ \hline
Binary response is meaningful & Undecided voters or measurement error & Richer outcome/measurement model \\ \hline
Prior relevance & Historical regime changed & Skeptical priors; robustness analysis \\
\bottomrule
\end{tabular}
\end{center}
\end{frame}

%-----------------------------------------------
\begin{frame}{Ethics belongs inside the risk register}
Before collecting, modeling, or acting, ask:
\begin{itemize}
\item Who receives benefits, and who bears harms?
\item Who is included, excluded, or made less visible by the sample?
\item Are consent, privacy, and data-use expectations respected?
\item What harms arise when results are generalized beyond the study population?
\item Were affected communities involved in defining the decision and its losses?
\end{itemize}
\begin{block}{Professional responsibility}
Technical validity is necessary,\\but it does not establish that a decision is legitimate or ethical.
\end{block}
\end{frame}

{
\setbeamercolor{background canvas}{bg=orange!10}
%-----------------------------------------------
\begin{frame}{In-class Exercise}
An inspection system has prior defect rate
\[
\pi\sim\mathrm{Beta}(2,38).
\]
A new audit finds $y=4$ defective items among $n=20$.
\begin{enumerate}
\item Write the posterior distribution.
\item Compute its mean.
\item Define an action threshold using false-alarm cost $C_{FP}=1$ and missed-defect cost $C_{FN}=9$.
\item What posterior probability must be estimated to decide whether to intervene when ``high risk'' means $\pi>0.10$?
\item Name one model-risk concern that the calculation omits.
\end{enumerate}
\end{frame}

%-----------------------------------------------
\begin{frame}{}
\textbf{Answer:}
{\color{blue}
\[
\pi\mid y\sim\mathrm{Beta}(2+4,38+16)=\mathrm{Beta}(6,54),
\qquad \E(\pi\mid y)=0.10.
\]
Let $q=\Prb(\pi>0.10\mid y)$. Intervene when
\[
q>\frac{C_{FP}}{C_{FP}+C_{FN}}=\frac{1}{10}=0.10
\]
Because the posterior mean equals the risk threshold, $q$ is near one-half, so intervention is likely under this asymmetric loss.

\vspace{0.5em}
Possible model risks: clustered defects, changing production conditions, imperfect detection, or a nonrepresentative audit.
}
\end{frame}
}

%-----------------------------------------------
{
\setbeamercolor{background canvas}{bg=orange!10}
\begin{frame}{Practice Questions}
\small
\begin{enumerate}
\item If prior concentration doubles but its mean stays fixed, what changes?
\item Why is $L(\pi\mid y)$ not $\Prb(\pi\mid y)$?
\item What does conjugacy buy us computationally?
\item When can two analysts share a posterior but rationally choose different actions?
\item Which assumption would you challenge first in the election-poll example?
\item Identify the most decision-sensitive assumption.
\item Provide one example where false-positive and false-negative costs differ.
% \item Questions in the textbook: Exercise 3.4, Exercise 3.9, Exercise 3.11.
\end{enumerate}
\end{frame}

\setbeamercolor{background canvas}{bg=orange!10}
\begin{frame}{Practice Questions}
\small
\begin{enumerate}
\item[8.] Six Beta pdfs are plotted below.\\
Match each to one of the following models:\\
Beta(0.5,0.5), Beta(1,1), Beta(2,2), Beta(6,6), Beta(6,2), Beta(0.5,6).\\
\begin{center}
\includegraphics[height=4cm]{unnamed-chunk-114-1.png}
\end{center}
\end{enumerate}
\end{frame}

%-----------------------------------------------
\begin{frame}{}
\small
\begin{enumerate}
\item[9.] What do you call a sweet carbonated drink: pop, soda, coke, or something else?
Let $\pi$ be the proportion of U.S. residents that prefer the term "pop." Two different beverage salespeople from different regions of the country have different priors for $\pi$.
The first salesperson works in North Dakota and specifies a Beta(8,2) prior. The second works in Louisiana and specifies a Beta(1,20) prior.
\begin{enumerate}
\item[a. ] Calculate the prior mean, mode, standard deviation of $\pi$ for both salespeople.
\item[b. ] Plot the prior pdfs for both salespeople.
\item[c. ] Compare, in words, the salespeople’s prior understandings about the proportion of U.S. residents that say “pop."
\end{enumerate}
\end{enumerate}
\end{frame}

%-----------------------------------------------
\begin{frame}{}
\small
\begin{enumerate}
\item[10.] A university wants to know what proportion of students are regular bike riders, $\pi$, so that they can install an appropriate number of bike racks.
Since the university is in sunny Texas, staff think that $\pi$ has a mean of 1 in 4 students, and a mode of $5/22$.
\begin{enumerate}
\item[a. ] Specify and plot a Beta model that reflects the staff’s prior ideas about $\pi$.
\item[b. ] Among 50 surveyed students, 15 are regular bike riders. What is the posterior model for $\pi$?
\item[c. ] What is the mean, mode, and standard deviation of the posterior model?
\item[d. ] Does the posterior model more closely reflect the prior information or the data? Explain your reasoning.
\end{enumerate}
\end{enumerate}
\end{frame}
}

%-----------------------------------------------
\begin{frame}{The model is simple; disciplined use is not}
\begin{enumerate}
\item Specify $\pi$, the action, and the consequence-relevant threshold.
\item Defend the Beta prior's center and concentration.
\item Check whether the Binomial assumptions fit the data-generating process.
\item Update: $\mathrm{Beta}(\alpha+y,\beta+n-y)$.
\item Compute posterior and predictive quantities tied to the decision.
\item Minimize posterior expected loss.
\item Stress-test prior, likelihood, losses, and ethical impacts.
\end{enumerate}
\begin{block}{}
Bayesian updating quantifies learning; decision analysis makes the consequences explicit.
\end{block}
\end{frame}

\end{document}